\documentclass[../../main.tex]{subfiles}

\begin{document}

\section{充满空间的曲线}\label{section:ch07-space-filling-curve}

当 $Y$ 是完备度量空间时, $\mathcal{C}(X,Y)$ 在一致度量下是完备的. 作为这一完备性的应用, 我们将要构造著名的"充满空间的 Peano 曲线".

\begin{theorem}\label{theorem:ch07-peano-curve}
设 $I=[0,1]$. 则存在一个连续映射 $f\colon I\to I^2$, 它的像填满整个正方形 $I^2$.
\end{theorem}

这种道路的存在性, 如同无处可微的连续函数 (后面将要讲到) 的存在性一样, 违反了人们朴素的集合直觉.

\par
\begin{figure}[htbp]
\centering
\includegraphics[width=0.85\textwidth]{fig44-1.png}
\caption{}\label{figure:ch07-fig44-1}
\end{figure}

\par
\begin{figure}[htbp]
\centering
\includegraphics[width=0.85\textwidth]{fig44-2.png}
\caption{}\label{figure:ch07-fig44-2}
\end{figure}

\par
\begin{figure}[htbp]
\centering
\includegraphics[width=0.85\textwidth]{fig44-3.png}
\caption{}\label{figure:ch07-fig44-3}
\end{figure}

\par
\begin{figure}[htbp]
\centering
\includegraphics[width=0.7\textwidth]{fig44-4.png}
\caption{}\label{figure:ch07-fig44-4}
\end{figure}

\par
\begin{figure}[htbp]
\centering
\includegraphics[width=0.7\textwidth]{fig44-5.png}
\caption{}\label{figure:ch07-fig44-5}
\end{figure}

\begin{proof}
第一步. 我们将构造一个映射 $f$, 使得它是一个连续函数序列 $f_n$ 的极限. 为了给出序列 $f_n$, 我们首先来描述关于道路的一种特殊运算.

我们从实直线上任意闭区间 $[a,b]$ 和平面上其边平行于坐标轴的任意正方形入手, 考虑如\cref{figure:ch07-fig44-1}所示的三角形道路 $g$, 它是 $[a,b]$ 到这个正方形的一个连续映射. 我们要描述的运算就是用\cref{figure:ch07-fig44-2}所示的道路 $g'$ 代替道路 $g$, 它由四条三角形道路组成, 其中的每一个只有 $g$ 的一半那么大. 注意 $g$ 与 $g'$ 有相同的起点和终点. 如果愿意的话, 你可以写出 $g$ 和 $g'$ 的方程.



对于连接正方形两个相邻顶点的任意三角形道路, 都可以施行上面所说的运算. 例如, 对于\cref{figure:ch07-fig44-3}所示的道路 $h$, 重复上述运算就得出道路 $h'$.


第二步. 现在定义函数序列 $f_n\colon I\to I^2$. 为方便起见, 第一个函数 $f_0$ 就是\cref{figure:ch07-fig44-1}所示的三角形道路, 其中 $a=0$, $b=1$. 下一个函数 $f_1$ 是对函数 $f_0$ 施行第一步中所述运算而得到的函数, 如\cref{figure:ch07-fig44-2}所示. 第三个函数 $f_2$ 是对组成 $f_1$ 的四条三角形道路中的每一个施行同样的运算而得出的函数, 如\cref{figure:ch07-fig44-4}所示.


第四个函数 $f_3$ 则是对组成 $f_2$ 的十六条三角形道路中的每一个施行上述运算而得到的函数, 如\cref{figure:ch07-fig44-5}所示. 依此类推. 一般地, $f_n$ 是由在第一步中考虑过的 $4^n$ 条三角形道路组成的道路, 其中的每一个都落在边长为 $1/2^n$ 的一个正方形中. 对这些三角形道路施行第一步所描述的运算, 每一个用四条较小的三角形道路来代替, 就得到函数 $f_{n+1}$.


第三步. 出于证明的需要, 用 $d(\mathbf{x},\mathbf{y})$ 表示 $\mathbb{R}^2$ 中的平方度量,
\begin{align*}
d(\mathbf{x},\mathbf{y})=\max\{|x_1-y_1|,|x_2-y_2|\}.
\end{align*}
于是, 我们令 $\rho$ 表示相应的 $\mathcal{C}(I,I^2)$ 上的上确界度量,
\begin{align*}
\rho(f,g)=\sup\{d(f(t),g(t))\mid t\in I\}.
\end{align*}
因为 $I^2$ 在 $\mathbb{R}^2$ 中是闭集, 所以它在平方度量下是完备的, 于是 $\mathcal{C}(I,I^2)$ 关于度量 $\rho$ 是完备的.

我们断言, 第二步所定义的函数序列 $(f_n)$ 是度量 $\rho$ 下的一个 Cauchy 序列. 为此, 我们考察从 $f_n$ 到 $f_{n+1}$ 时会发生什么变化. 构成 $f_n$ 的每一个小三角形道路都落在边长为 $1/2^n$ 的某一个正方形中. 而从 $f_n$ 得出 $f_{n+1}$ 的运算就是用同一正方形中的四条三角形道路来代替一条三角形道路. 因此, 在 $I^2$ 的平方度量下, $f_n(t)$ 和 $f_{n+1}(t)$ 之间的距离最多为 $1/2^n$. 于是有 $\rho(f_n,f_{n+1})\leqslant 1/2^n$. 由此推出 $(f_n)$ 是一个 Cauchy 序列, 因为对于所有的 $n$ 和 $m$, 有
\begin{align*}
\rho(f_n,f_{n+m})\leqslant 1/2^n+1/2^{n+1}+\cdots+1/2^{n+m-1}<2/2^n.
\end{align*}

第四步. 因为 $\mathcal{C}(I,I^2)$ 是完备的, 所以序列 $f_n$ 收敛于一个连续函数 $f\colon I\to I^2$. 下面我们证明 $f$ 是一个满射.

设 $\mathbf{x}$ 是 $I^2$ 的一个点, 我们证明 $\mathbf{x}\in f(I)$. 首先, 我们注意, 对于给定的 $n$, 道路 $f_n$ 能够到达距离点 $\mathbf{x}$ 为 $1/2^n$ 的范围以内, 因为把 $I^2$ 分割为边长是 $1/2^n$ 的小正方形, 道路 $f_n$ 会碰到每一个这样的小方块.

利用这一点, 我们证明, 对任意 $\varepsilon>0$, $\mathbf{x}$ 的 $\varepsilon$-邻域都与 $f(I)$ 相交. 选取充分大的 $N$, 使得
\begin{align*}
\rho(f_N,f)<\varepsilon/2\quad\text{且}\quad 1/2^N<\varepsilon/2.
\end{align*}
根据前一段的结果, 存在一个点 $t_0\in I$, 使得 $d(\mathbf{x},f_N(t_0))\leqslant 1/2^N$. 因为对于所有的 $t$, 有 $d(f_N(t),f(t))<\varepsilon/2$, 所以
\begin{align*}
d(\mathbf{x},f(t_0))<\varepsilon,
\end{align*}
从而 $\mathbf{x}$ 的 $\varepsilon$-邻域与 $f(I)$ 相交.

由此可见 $\mathbf{x}$ 属于 $f(I)$ 的闭包. 但由于 $I$ 是紧致的, 所以 $f(I)$ 是紧致的, 从而是闭的. 于是证明了 $\mathbf{x}$ 属于 $f(I)$.
\end{proof}

\subsection*{习题}

\begin{enumerate}
\item 给定 $n$, 证明存在一个连续的满射 $g\colon I\to I^n$. [提示: 考虑 $f\times f\colon I\times I\to I^2\times I^2$.]
\item 证明存在连续满射 $f\colon\mathbb{R}\to\mathbb{R}^n$.
\item 
\begin{enumerate}[(a)]
\item 若 $\mathbb{R}^\omega$ 被赋予积拓扑, 证明不存在连续的满射 $f\colon\mathbb{R}\to\mathbb{R}^\omega$. [提示: 证明 $\mathbb{R}^\omega$ 不能表示成可数个紧致子空间的并.]
\item 若 $\mathbb{R}^\omega$ 被赋予积拓扑, 确定是否存在连续满射 $f\colon\mathbb{R}\to\mathbb{R}^\infty$.
\item 若 $\mathbb{R}^\omega$ 被赋予一致拓扑或箱拓扑, 对于小题 (a) 和 (b) 会有怎样的结论呢?
\end{enumerate}
\item
\begin{enumerate}[(a)]
\item 设 $X$ 是一个 Hausdorff 空间. 证明: 若存在连续满射 $f\colon I\to X$, 则 $X$ 是紧致的、连通的、弱局部连通的, 并且可以度量化. [提示: 证明 $f$ 是完备映射.]
\item 上面小题 (a) 的逆命题也成立, 它是点集拓扑学中著名的 Hahn-Mazurkiewicz 定理 (见 [H-Y] p. 129). 假设我们已知这个定理, 然后证明存在一个连续满射 $f\colon I\to I^\omega$.
\end{enumerate}
\end{enumerate}

如果某一个 Hausdorff 空间是单位闭区间的连续像, 通常称之为一个 Peano 空间(Peano space).

\end{document}